\chapter*{General Conclusion}
\addcontentsline{toc}{chapter}{General Conclusion}

This report serves as a comprehensive study and implementation of an end-to-end Edge-AI computer vision system, developed to run natively on Neural Processing Units (NPUs) for the EYE-D video analytics solution. Our work strategy was inspired by the CRISP-DM methodology, ensuring a structured and iterative approach to overcoming complex industrial challenges.

We began by understanding the business context and the strict hardware requirements, specifically analyzing the severe SRAM memory and compute constraints inherent to edge devices. This allowed us to identify the critical limitations in traditional cloud-centric visual processing workflows and highlight the necessity of migrating inference to the edge. Following this, we established a rigorous data engineering lifecycle, utilizing advanced annotation and analytics tools to curate highly representative datasets for both logistical tracking and industrial safety domains. 

Subsequently, we selected, fine-tuned, and optimized single-stage object detection architectures to achieve a precise balance between high-accuracy localization and real-time inference speeds. We addressed critical environmental artifacts and challenging visual conditions, ultimately converting and quantizing the models into NPU-compatible formats. Furthermore, we explored advanced Vision-Language Models and Zero-Shot detection paradigms to evaluate their viability for future, highly adaptable edge deployments.

The achieved results demonstrate the immense potential of Edge-AI in real-world industrial and safety environments. By embedding intelligent processing directly at the edge, our solution ensures rapid responsiveness, strict data privacy, and robust performance, providing a scalable and highly modular foundation for future intelligent video analytics.
